\documentclass[../main]{subfiles}
\ifSubfilesClassLoaded{\setcounter{chapter}{12}}{}
\begin{document}

\chapter{Pembuktian Tidak Langsung dan Kontradiksi}

\begin{subcpmk}
  \item \textbf{Sub-CPMK 4.1:} Mahasiswa mampu menerapkan berbagai metode pembuktian matematis dasar (pembuktian langsung, pembuktian tidak langsung/kontradiksi, dan kontraposisi).
\end{subcpmk}

% ============================================================
% MATERI POKOK
% ============================================================
\section[Pembuktian dengan Kasus]{\raggedright Pembuktian dengan Kasus dan Pembuktian Exhaustif}

Terdapat situasi di mana membuktikan suatu pernyataan secara umum (untuk seluruh elemen domain) lebih mudah dilakukan dengan membagi domain menjadi beberapa kasus yang saling lepas (\textit{mutually exclusive}) dan menangkap seluruh kemungkinan (\textit{collectively exhaustive}), kemudian membuktikan pernyataan tersebut untuk setiap kasus secara terpisah \cite{rosen2019}.

\subsection{Pembuktian dengan Kasus (Proof by Cases)}

\begin{definisi}[Pembuktian dengan Kasus]
\logic{Pembuktian dengan kasus} (\textit{proof by cases} / \textit{proof by exhaustion of cases}) adalah teknik pembuktian di mana domain dipartisi menjadi kasus-kasus $C_1, C_2, \ldots, C_n$ yang saling lepas dan lengkap, kemudian pernyataan dibuktikan secara terpisah untuk setiap kasus. Teknik ini valid berdasarkan aturan inferensi:
\[ \frac{(C_1 \rightarrow q) \wedge (C_2 \rightarrow q) \wedge \ldots \wedge (C_n \rightarrow q)}{(C_1 \vee C_2 \vee \ldots \vee C_n) \rightarrow q} \]
\end{definisi}

Kunci keberhasilan metode ini adalah memastikan bahwa kasus-kasus yang dipilih benar-benar \emph{mencakup seluruh kemungkinan}. Jika ada kasus yang terlewat, bukti menjadi tidak valid.

\begin{contoh}
\textbf{Teorema}: Untuk setiap bilangan bulat $n$, $n^2 + n$ adalah bilangan genap.

\textbf{Bukti (dengan Kasus):}

Setiap bilangan bulat $n$ bernilai genap atau ganjil (dua kasus, saling lepas dan lengkap).

\textbf{Kasus 1:} $n$ genap. Maka $n = 2k$ untuk suatu bilangan bulat $k$.
\[ n^2 + n = (2k)^2 + 2k = 4k^2 + 2k = 2(2k^2 + k) \]
yang berbentuk $2m$ (genap). $\checkmark$

\textbf{Kasus 2:} $n$ ganjil. Maka $n = 2k + 1$ untuk suatu bilangan bulat $k$.
\[ n^2 + n = (2k+1)^2 + (2k+1) = 4k^2 + 4k + 1 + 2k + 1 = 4k^2 + 6k + 2 = 2(2k^2 + 3k + 1) \]
yang berbentuk $2m$ (genap). $\checkmark$

Kedua kasus terbukti. $\blacksquare$
\end{contoh}

\subsection{Pembuktian Exhaustif}

\begin{definisi}[Pembuktian Exhaustif]
\logic{Pembuktian exhaustif} (\textit{exhaustive proof} / \textit{brute-force proof}) adalah pembuktian di mana pernyataan diverifikasi untuk \textbf{setiap} elemen domain secara individual, tanpa menggunakan argumentasi umum. Metode ini hanya dapat diterapkan jika domain bersifat berhingga dan cukup kecil untuk diperiksa satu per satu.
\end{definisi}

\begin{contoh}
\textbf{Teorema}: Untuk setiap bilangan bulat $n$ dengan $1 \leq n \leq 4$, berlaku $n^2 \leq 2^n \cdot n$.

\textbf{Bukti (Exhaustif):} Periksa setiap nilai:
\begin{center}
\renewcommand{\arraystretch}{1.2}
\begin{tabular}{|c|c|c|c|}
\hline
$n$ & $n^2$ & $2^n \cdot n$ & $n^2 \leq 2^n \cdot n$? \\ \hline
1 & 1 & 2 & Ya ($1 \leq 2$) \\ \hline
2 & 4 & 8 & Ya ($4 \leq 8$) \\ \hline
3 & 9 & 24 & Ya ($9 \leq 24$) \\ \hline
4 & 16 & 64 & Ya ($16 \leq 64$) \\ \hline
\end{tabular}
\end{center}

Semua kasus terverifikasi. $\blacksquare$
\end{contoh}

\begin{catatan}
Pembuktian exhaustif memiliki keterbatasan yang jelas: jika domain berhingga tetapi sangat besar (misalnya semua bilangan prima kurang dari $10^9$), verifikasi manual menjadi tidak praktis. Dalam kasus seperti itu, komputer sering digunakan untuk membantu verifikasi, tetapi bukti semacam itu kadang diperdebatkan statusnya dalam komunitas matematika.
\end{catatan}


% ============================================================
% AKTIVITAS PEMBELAJARAN
% ============================================================
\begin{aktivitas}
  \item Buktikan secara kontradiksi: ``$\sqrt{3}$ adalah bilangan irasional.'' (Petunjuk: ikuti pola bukti irasionalitas $\sqrt{2}$.)
  \item Buktikan dengan kontradiksi: ``Jika $n^2$ habis dibagi 3, maka $n$ habis dibagi 3.''
  \item Diskusikan perbedaan antara bukti kontraposisi dan bukti kontradiksi. Kapan masing-masing lebih tepat digunakan?
\end{aktivitas}

% ============================================================
% LATIHAN DAN REFLEKSI
% ============================================================
\begin{latihan}
  \item Buktikan secara kontradiksi: ``Tidak ada bilangan bulat terbesar.''
  \item Buktikan: ``Jumlah bilangan rasional dan bilangan irasional adalah irasional'' (gunakan kontradiksi).
  \item \textbf{Refleksi}: Dalam konteks verifikasi perangkat lunak, bagaimana konsep kontradiksi (asumsi salah mengarah ke kemustahilan) terkait dengan teknik \textit{model checking}?
\end{latihan}

% ============================================================
% ASESMEN
% ============================================================
\begin{asesmen}
\textbf{Instrumen untuk Sub-CPMK 4.1 (Pembuktian Kontradiksi)}

\begin{enumerate}
  \item Pembuktian bahwa $\sqrt{2}$ irasional menggunakan metode:
  \begin{enumerate}
    \item Bukti langsung \item Bukti kontraposisi \item Bukti kontradiksi \item Bukti exhaustif
  \end{enumerate}
  Jawaban: (c)

  \item Buktikan dengan kontradiksi: ``Terdapat tak berhingga banyak bilangan prima.'' (Petunjuk: gunakan konstruksi $N = p_1 p_2 \cdots p_n + 1$.)
\end{enumerate}

\textbf{Rubrik}: Lihat Lampiran.
\end{asesmen}

% ============================================================
% CHECKLIST KOMPETENSI
% ============================================================
\begin{checklist}
  \item Saya mampu menyusun bukti kontradiksi dengan mengasumsikan negasi dan menurunkan kemustahilan
  \item Saya memahami perbedaan antara bukti kontraposisi dan bukti kontradiksi
  \item Saya dapat memilih penggunaan kontradiksi untuk pernyataan eksistensial negatif atau ketidakmungkinan
\end{checklist}

% ============================================================
% RANGKUMAN
% ============================================================
\begin{rangkuman}
Bab ini membahas \textbf{pembuktian tidak langsung dengan kontradiksi}. Untuk membuktikan pernyataan $P$, kita mengasumsikan $\neg P$ benar lalu menurunkan suatu kontradiksi (pernyataan yang mustahil benar). Karena asumsi mengarah ke kemustahilan, $P$ harus benar. Kontradiksi sangat berguna untuk membuktikan irasionalitas (misalnya $\sqrt{2}$), ketidakberhinggaan (misalnya bilangan prima), dan pernyataan negatif. Kontradiksi berbeda dari kontraposisi: kontraposisi hanya berlaku untuk implikasi $p \rightarrow q$ dengan membuktikan $\neg q \rightarrow \neg p$; kontradiksi berlaku untuk semua jenis pernyataan.
\end{rangkuman}

\ifSubfilesClassLoaded{
  \renewcommand{\bibname}{Daftar Pustaka}
  \bibliographystyle{plain}
  \bibliography{../references}
}{}
\end{document}
